\section{Теория на корутините и асинхронното изпълнение в C++}
\label{sec:coroutines_async_theory}

Корутините са езиков механизъм за описване на прекъсваеми изчисления, при които изпълнението може да бъде спряно и по-късно възобновено без загуба на локалното състояние. В стандарта C++20 корутините влизат в основния език, а не се реализират като външна библиотечна симулация~\cite{cpp20standard,p0912r5,p1063r0}. Това е важно за настоящата работа, защото позволява асинхронният слой на библиотеката да се изгражда върху дефинирани по време на компилация правила за типове, жизнен цикъл и поток на управление (поток на управление), вместо върху неструктурирана архитектура от функции за обратно извикване (функция за обратно извикване).

В тази глава корутините се разглеждат не като удобен синтаксис, а като изпълнителен модел, който стои между чисто синхронното блокиращо изпълнение и тежките нишки на операционната система. Целта не е да се покаже само как се използват \code{co_await} и \code{co_return}, а да се изведе защо преобразуването от компилатора към автомат с крайни състояния прави този подход особено подходящ за библиотека, която и без това се основава на статични типове, специализации на характеристики и валидиране по време на компилация.

Този теоретичен пласт е необходим и поради още една причина. В контекста на библиотеката за обектно-релационно съпоставяне по време на компилация корутините не са изолиран езиков експеримент, а основа за единен асинхронен интерфейс над свързващи компоненти с много различни експлоатационни характеристики. За да бъде такава архитектура научно защитима, трябва ясно да се разграничат езиковият механизъм, инфраструктурата по време на изпълнение и конкретната семантика на изпълнение на драйверите.

\subsection{Процеси, нишки, задачи и корутини}

Процесът е защитен адресен контекст на операционната система, нишката е планирана от ядрото последователност от инструкции, а задачата е по-общо програмно понятие за единица работа. Корутината не е нито процес, нито нишка. Тя е кооперативна единица за изпълнение, чието спиране и подновяване се управлява от програмния модел и генерираната структура по време на изпълнение, а не директно от планировчика на операционната система.

Това разграничение е съществено, защото библиотеката комбинира корутини и набор от нишки (\code{thread\_pool}), без да приравнява двете абстракции. Нишката е скъпа ресурсна единица на ядрото, свързана с отделен стек, планиране и примитиви за синхронизация. Корутината е значително по-лека логическа единица, която пази само състоянието, необходимо за спиране и продължаване на конкретно изчисление. Следователно преминаването към модел, базиран на корутини, не означава „повече нишки“, а по-структуриран начин за описване на това къде едно изчисление може да изчака външно събитие.

\begin{table}[H]
\centering
\caption{Съпоставка между основните изпълнителни единици}
\label{tab:execution_models_comparison}
\begin{tabularx}{\textwidth}{L{2.8cm}L{4.0cm}Y}
\toprule
\textbf{Единица} & \textbf{Същност} & \textbf{Значение за библиотеката} \\
\midrule
Процес & Изолиран адресен контекст с ресурси на операционната система & Не е пряка единица на архитектурата, но определя границата на средата за изпълнение \\
Нишка & Планиран от ядрото поток на изпълнение със собствен стек & Използва се от \code{thread\_pool} за изнасяне (изнасяне) на блокиращи операции \\
Корутина & Кооперативно прекъсваемо изчисление със запазено локално състояние & Осигурява формалния модел за \code{Task<T>} и за асинхронната композиция \\
\bottomrule
\end{tabularx}
\end{table}

Таблица~\ref{tab:execution_models_comparison} подчертава, че корутината е логическа, а не единица на ядрото. Именно това позволява в една и съща нишка да се редуват множество спиране/възобновяване последователности, без всяка от тях да изисква отделен стек и отделна нишкова идентичност.

\subsection{Синтаксис, протокол за изчакване и езикови свързванеing-и}

Ключовите езикови оператори \code{co_await}, \code{co_return} и \code{co_yield} маркират, че дадена функция следва семантиката на корутините. От гледна точка на архитектурата най-важен е \code{co_await}, защото именно той описва точката, в която текущото изчисление може да отстъпи изпълнението, докато външна операция --- например база данни или планиране към набор от нишки --- не стане готова. Така синтаксисът на потребителя остава близък до линейния код, докато изпълнителният модел става асинхронен.

За текущата дипломна работа е важно, че тази езикова форма не е динамично „откривана“ по време на изпълнение, а е компилаторно свързана с типа на връщаната стойност на функцията и с \code{promise\_type}, който \code{std::coroutine\_traits} извежда за конкретния подпис~\cite{cpp20standard,p0912r5}. Това означава, че начинът, по който една корутина създава резултат, предава грешка или съхранява продължение (продължение), може да бъде формализиран още на ниво типова система.

На практика \code{co_await} работи чрез протокол за изчакване (протокол за изчакване). Обектът, който се изчаква, участва в тристъпкова схема: \code{await\_ready()} определя дали спиране изобщо е необходимо, \code{await\_suspend()} получава манипулатор на корутината (манипулатор на корутината) и решава как да организира възобновяването, а \code{await\_resume()} връща резултата или повдига грешка. Точно тази тройка прави възможно различни източници на асинхронност --- задачи в набор от нишки, готовност на цикъл за събития, функции за обратно извикване от драйвер --- да бъдат уеднаквени под общ \code{co_await} синтаксис.

\subsection{Корутината като автомат с крайни състояния}

Компилаторът не „изпълнява магия“, а преобразува корутинната функция в автомат с крайни състояния с ясно определени точки на спиране, преходи на възобновяване и съхранение на рамки~\cite{p0912r5,p1063r0}. Това преобразуване е от особено значение за настоящата теза, защото показва, че подходът по време на компилация на библиотеката не се изчерпва само с моделиране на заявки. Същият езиков инструментариум позволява и асинхронният поток на управление да бъде формализиран чрез типове и статични правила, вместо да се описва чрез неструктурирани вериги от функции за обратно извикване.

Концептуално една корутина функция преминава през няколко ясно разграничени фази: създаване на рамка и \code{promise\_object}, начално изпълнение, достигане до точка на спиране, външно организирано възобновяване и финално завършване. Това е съществено, защото коректност вече не зависи само от „какво прави кодът“, а и от това кога е допустимо той да бъде продължен и кой носи отговорност за продължаването.

\begin{figure}[H]
\centering
\begin{tikzpicture}[node distance=1.5cm and 1.8cm]
    \node[ormblock, text width=3.1cm, align=center] (call) {Извикване на корутинна функция};
    \node[ormblock, text width=3.4cm, align=center, right=of call] (рамка) {Създаване на корутинна рамка и \code{promise\_object}};
    \node[ormblock, text width=2.8cm, align=center, right=of рамка] (run) {Начално изпълнение на тялото};

    \node[ormaccent, text width=2.9cm, align=center, below=of рамка] (await) {\code{co\_await} върху обект};
    \node[ormblock, text width=3.0cm, align=center, right=of await] (спиране) {Състояние на спряна корутина};
    \node[ormblock, text width=3.2cm, align=center, below=of спиране] (source) {Източник на възобновяване:\\набор от нишки / цикъл за събития / драйверна функция за обратно извикване};
    \node[ormaccent, text width=2.8cm, align=center, left=of source] (възобновяване) {Възобновяване чрез манипулатор};

    \node[ormblock, text width=2.9cm, align=center, left=of възобновяване] (finish) {Продължаване на тялото и \code{co\_return}};
    \node[ormblock, text width=3.0cm, align=center, below=of finish] (final) {\code{final\_suspend}, продължение, резултат/грешка};

    \coordinate (run_спиране_mid) at ($(run)!0.5!(спиране)$);

    \draw[ormline] (call) -- (рамка);
    \draw[ormline] (рамка) -- (run);
    \draw[ormline] (run) -- (run |- run_спиране_mid) -- (await |- run_спиране_mid) -- (await.north);
    \draw[ormline] (await) -- (спиране);
    \draw[ormline] (спиране) -- (source);
    \draw[ormline] (source) -- (възобновяване);
    \draw[ormline] (възобновяване) -- (finish);
    \draw[ormline] (finish) -- (final);
\end{tikzpicture}
\caption{Концептуален жизнен цикъл на корутината като рамка, точки на спиране и външно управлявано възобновяване}
\label{fig:coroutine_lifecycle}
\end{figure}

Фигура~\ref{fig:coroutine_lifecycle} показва защо корутинният модел е естествен за библиотека с ясно отделени слоеве. Логическият код на потребителя остава линеен, но реалното управление на изпълнението преминава през формални преходи, които могат да бъдат свързани с конкретни механизми по време на изпълнение.

\subsection{Корутинна рамка, обект-обещание, продължение и безопасност}

При всяка корутина компилаторът създава корутинна рамка, в която се съхраняват локалните променливи, текущото състояние на изпълнението и обектът-обещание (обект-обещание), чрез който корутината комуникира резултат, грешка или завършване. Това има директно отражение върху безопасността: жизненият цикъл на рамката, собствеността върху резултата и мястото на продължението не са второстепенни детайли, а част от договора за коректност на асинхронния слой.

Трябва да се подчертае, че корутинната рамка не е обикновена стекова рамка. Тя може да преживее излизането от текущия стеков контекст и да бъде възобновена по-късно от друга нишка или от инфраструктура за цикъл за събития. Поради това проблеми като преждевременно разрушаване на рамката, двойно възобновяване, неправилно предаване на състояние за изключение или загуба на продължението не са „детайли на реализацията на ниско ниво“, а фундаментални рискове за коректността.

Именно обектът-обещание организира връщането на стойност, разпространението на грешка и поведението при \code{initial\_suspend} и \code{final\_suspend}. От гледна точка на архитектурата това е важен пример за свързване по време на компилация: видът на резултата и семантиката на завършването не се определят от динамична схема за регистрация, а са закодирани в типа на връщаната стойност на корутината. Това позволява \code{Task<T>} да бъде проектиран като формална абстракция, а не като неформален обект, подобен на бъдеще (бъдеще).

\subsection{Корутини, набор от нишки и инфраструктура за цикъл за събития}

Наборът от нишки (\code{thread\_pool}) и корутините решават различни проблеми. Наборът от нишки предоставя ограничен брой работни нишки и политика за планиране на работа, докато корутините описват как една логическа задача може да бъде спряна и по-късно възобновена. Комбинацията между двете позволява блокиращи свързващи компоненти да бъдат интегрирани по универсален начин чрез изнасяне (изнасяне), без потребителят да се връща към програмиране, ориентирано към функции за обратно извикване.

От друга страна, при истински неблокиращи драйвери корутината може да бъде възобновявана не от работна нишка, а от цикъл за събития или мост за функции за обратно извикване. Средоподобни на реактори наблюдават готовността на файлови дескриптори чрез системни механизми като \code{kqueue} на Mac OS и BSD системите~\cite{kqueue_man}. Следователно корутината сама по себе си не „изпълнява входно-изходни операции“, а предоставя формален интерфейс, през който готовността за вход-изход или функцията за обратно извикване на драйвера може да се превърне в действие за възобновяване.

Точно тук става ясно защо корутините не трябва да бъдат смесвани с нишките. Нишката е носител на физическо изпълнение; корутината е носител на логическа последователност. Един и същ асинхронен модел може да използва набор от нишки в един сценарий на свързващ компонент и цикъл за събития в друг, без да променя потребителския синтаксис на \code{co_await}.

\subsection{Значение за библиотеката за обектно-релационно съпоставяне по време на компилация}

За разработваната библиотека този теоретичен анализ има пряко архитектурно следствие. Моделираното по време на компилация междинно представяне на заявкитете и договорът на конектора определят какво е позволено да се изпълни; базираният на корутини слой определя как това позволено изчисление може да бъде изпълнено асинхронно. С други думи, типово коректната заявка и коректният път на изпълнение са различни, но взаимно зависими пластове.

Това разграничение е особено важно при библиотека с множество свързващи компоненти. PostgreSQL и MySQL предоставят платформено-специфични неблокиращи клиентски интерфейси, но с различни експлоатационни форми~\cite{libpq_async,mysql_capi_async}. \code{hiredis} използва асинхронен слой, базиран на функции за обратно извикване~\cite{hiredis}, а драйверът за Cassandra работи чрез механизъм с бъдеща стойност и функция за обратно извикване~\cite{cassandra_futures}. При MongoDB фокусът остава върху безопасен пул и многонишкова употреба на \code{libmongoc}~\cite{mongoc_client_pool}, докато \code{libneo4j-client} следва преобладаващо синхронен модел заявка/отговор~\cite{libneo4j_client}. Именно тази нееднородност оправдава асинхронна архитектура с две стратегии: универсален модел за изнасяне (изнасяне) и платформено-специфична асинхронна интеграция за драйвери, които я позволяват.

Следователно корутините не правят автоматично всяка операция неблокираща и не отменят ограниченията на даден драйвер. Те дават формален модел за асинхронно описание на изчислението; реалният ефект върху производителността зависи от това дали клиентската библиотека има платформено-специфичен неблокиращ интерфейс, дали работата се изнася към набор от нишки или дали задачата остава синхронна. Именно затова следващите глави разглеждат отделно ОРС архитектурата по време на компилация, асинхронния слой на рамката и семантиката на изпълнение на конкретните свързващи компоненти.
